\section{Solvability and rank logic}\label{sec:linear-algebraic-logics}

We now prove the linear-algebraic consequence stated in the introduction.
The query $Q$ is defined on all finite structures. For
Corollary~\ref{cor:linear-algebraic-logics} we take a nonempty input, so
the matrix index sets below are nonempty. We use the two-sorted syntax
of $\mathrm{FPS}_3$ and $\mathrm{FPR}_3$ in
\cite[Section~2]{graedelPakusa2019}.
We first define its matrix without ordering the input, then encode
consistency using the respective operators.

\begin{proof}[Proof of Corollary~\ref{cor:linear-algebraic-logics}]
Write $X(a)$ and $Y(a)$ for the predicates defining $X$ and $Y$, and put
$n=\#x(x=x)$, a closed numeric term. For $\delta=1,2$, the counting terms
\[
 d_\delta(y,a)=\#x\bigl(Y(x)\wedge\mathsf I(a,x)
                         \wedge\mathsf Z_\delta(y,x)\bigr)
\]
give, for each fixed $r\in\{0,1,2\}$, the FPC predicate
\[
 C_r(y,a)\ \equiv\
 \exists\eta\leq 3n\;\bigl(d_1(y,a)+2d_2(y,a)=3\eta+r\bigr).
\]
Here $\eta$ is a number variable. Since $d_1+2d_2\leq3n$, this says
exactly that $C(y,a)=r$ in $\mathbb F_3$. It includes both contributions
when $\mathsf Z_1$ and $\mathsf Z_2$ both hold, while $\mathsf Z_0$
contributes zero. The row test is the first-order predicate
\[
\begin{split}
 T(t,y)\ \equiv\ X(t)\wedge Y(y)\wedge\exists a\,\exists x\,
 \bigl(&X(a)\wedge\mathsf{VB}(t,a)\wedge Y(x)\\
       &{}\wedge\mathsf I(a,x)\wedge\mathsf{EB}(y,x)\bigr).
\end{split}
\]

For nonempty $A$, use the full tuple domain $D=A^2\times\{0,1\}$ for both
rows and columns, with the last coordinate in the number sort. For
$i=(t,y,\rho)$ and $j=(a,z,\kappa)$, define
\[
\begin{aligned}
 N(i)&\equiv \rho=0\wedge t=y\wedge X(t),&
 E(i)&\equiv \rho=1\wedge T(t,y),\\
 L(j)&\equiv \kappa=0\wedge a=z\wedge X(a),&
 U(j)&\equiv \kappa=1\wedge a=z\wedge Y(a).
\end{aligned}
\]
Set $b_i=\mathbf{1}_{N(i)}$ and, over $\mathbb F_3$, set
\[
 M_{ij}=
 \begin{cases}
 1&N(i)\wedge L(j)\wedge\mathsf{VB}(t,a),\\
 -C(y,a)&E(i)\wedge L(j)\wedge\mathsf{VB}(t,a),\\
 1&E(i)\wedge U(j)\wedge a=y,\\
 0&\text{otherwise}.
 \end{cases}
\]
The tags make these cases disjoint. The diagonal payloads give one column
for each $\lambda_a$ and one separately tagged column for each $\mu_y$,
even when $X$ and $Y$ overlap. They also give one normalization row for
each $t\in X$, while the $E$ rows retain every selected pair $(t,y)$.
All remaining rows and columns are zero, with right-hand side zero on
those rows. Thus $Mx=b$ is exactly the system defining $Q$, with harmless
zero padding. In particular an empty $X_t$ still gives $0=1$.
The predicates $B_r(i,j)\equiv(M_{ij}=r)$ for $r=1,2$ are FPC-definable
from these cases and $C_1,C_2$.

The solvability operator uses a Boolean matrix and an all-ones
right-hand side; compare the general fixed-ring normal form in
\cite[Lemma~4.1]{dghkp2013}. For an explicit encoding here, use full row
and column domains $D\times\{0,1,2,3\}$, with column tags denoting
$x^1_j,x^2_j,z_j,v_j$. For every $i\in D$, the four row tags encode
\[
\begin{gathered}
 \sum_{j:B_1(i,j)}x^1_j+
 \sum_{j:B_2(i,j)}(x^1_j+x^2_j)+\mathbf{1}_{\neg N(i)}v_i=1,\\
 v_i=1,\qquad x^1_i+z_i=1,\qquad x^2_i+z_i=1.
\end{gathered}
\]
All coefficients are zero or one: the two summands for $B_2$ use distinct
columns. Their incidence is one FPC formula using $B_1,B_2,N$, equality
of tuples in $D$, and the fixed numeric tags. Every row in the full
product has one of the four stated types. The last two equations force
$x^1=x^2$; together with $v=1$, the first equations then say $Mx^1=b$.
Conversely, a solution $x$ of $Mx=b$ extends by
$x^1=x^2=x$, $z=1-x$, and $v=1$. Hence one characteristic-three
solvability operator defines $Q$.

For rank, let $\mathbf{1}_\psi$ denote the numeric counting term
$\#\zeta\leq0.\psi$, where $\zeta$ is a number variable, and put
$m(i,j)=\mathbf{1}_{B_1(i,j)}+2\mathbf{1}_{B_2(i,j)}$.
On rows $D$ and columns $D\times\{0,1\}$, with numeric tag $h$, use the
entry terms
\[
\begin{aligned}
 \Theta_0(i;j,h)&=\mathbf{1}_{h=0}m(i,j),\\
 \Theta_+(i;j,h)&=\Theta_0(i;j,h)+\mathbf{1}_{h=1\wedge N(i)}.
\end{aligned}
\]
The first matrix is $M$ with zero columns; the second adjoins one copy
of $b$ for each $j\in D$. Since $D$ is nonempty, those identical copies
span the same space as one augmented column. The rank operator reduces
numeric entries modulo three, so
\[
 Q(A)\quad\Longleftrightarrow\quad
 \operatorname{rk}_3(\Theta_0)=\operatorname{rk}_3(\Theta_+).
\]
All domains used above are permitted Cartesian products with fixed
numeric bounds; none chooses or orders an input element.

Finally, the proof of Theorem~\ref{thm:main} uses only the nonempty
witnesses with $L\geq2$, so these definitions prove both noncontainments.
The empty input remains a true instance of $Q$ under its definition in
Section~\ref{sec:query}.
\end{proof}
